\chapter{Нейроны играют в теннис}
% полная версия: chapters/22_dishbrain.tex

Самый знаменитый опыт с живыми нейронами на чипе --- DishBrain, «мозг в чашке». Около миллиона нейронов, выращенных на чипе с двадцатью шестью тысячами электродов, играли в простейший компьютерный теннис: мяч летает по экрану, ракетка ходит вверх и вниз.

\section*{Как устроена игра}

Мяч попадает в сеть через восемь электродов. Какой из них щёлкает --- говорит, на какой высоте мяч. Как часто щёлкает --- говорит, далеко ли мяч: четыре раза в секунду, когда он у дальней стены, и всё чаще, до сорока, когда он подлетает к ракетке. Выход --- две области электродов: щелчки в одной двигают ракетку вверх, в другой --- вниз. Каждые десять миллисекунд компьютер считает щелчки и двигает ракетку. Эти десять миллисекунд --- такт компьютера стенда, не культуры: сама сеть работает асинхронно, как всякая нервная сеть.

\pencilfig{22}{\pencilart{22}}{Мяч, ракетка и миллион нейронов, которые её двигают.}

\section*{Учитель}

Никакого дофамина. После удачного удара сеть получает короткий и совершенно предсказуемый сигнал: все восемь электродов разом, одна десятая секунды. После промаха --- четыре секунды беспорядочной стимуляции в случайных местах и в случайные моменты, потом пауза, и новая подача.

\section*{Что вышло}

За двадцатиминутный сеанс розыгрыши у живых культур становились длиннее, а пропущенных подач --- меньше. Лучше всего это получалось при полной обратной связи. Когда вместо сигналов после удара и промаха просто наступала тишина, улучшение тоже было, но слабее; без обратной связи его не было. Человеческие нейроны играли немного лучше мышиных. А вот обучение, растянутое на несколько дней, надёжным не оказалось.

\section*{Почему сеть учится}

Авторы объясняют результат так: сеть стремится сделать свои ощущения предсказуемыми. Промах приносит хаос, удар --- порядок, и сеть меняется, чтобы порядка стало больше. Но есть и более приземлённое объяснение, которое тоже согласуется с данными: «встряхивай после промаха, не трогай после удара». Каждая встряска немного перемешивает связи наугад. Настройки, при которых сеть часто промахивается, всё время перемешиваются, а удачные остаются нетронутыми. Сеть сама собой задерживается там, где промахов мало, --- никто её не учил, её просто перестают трясти. В нашей простой модели такое правило действительно поднимает долю ударов.

\pencilfig{130}{%
% scrambler: miss probability p over one setting of the network (valley in the middle); a miss kicks the setting
% at random, a hit leaves it alone, so the time spent at a setting is proportional to 1/p (6x more in the valley here)
\draw[pen] (0,1.2) -- (8.2,1.2);
\draw[penlite=pcRed] plot[smooth] coordinates {(0.0,2.46) (0.8,2.46) (1.6,2.46) (2.0,2.44) (2.4,2.41) (2.8,2.32) (3.2,2.15) (3.6,1.90) (4.0,1.62) (4.4,1.44) (4.8,1.44) (5.2,1.62) (5.6,1.90) (6.0,2.15) (6.4,2.32) (6.8,2.41) (7.2,2.44) (7.6,2.46) (8.0,2.46)};
\node[lbl, text=pcRed, anchor=south] at (7.55,2.55) {промахов много};
\node[lbl, text=pcRed, anchor=west] at (5.3,1.45) {мало};
% kicked balls on the slopes
\draw[dotpen=pcBlue, wash=pcBlue] (1.3,2.68) circle (0.12);
\draw[arr=pcGray, densely dashed] (1.35,2.82) to[bend left=50] (2.35,2.72);
\draw[arr=pcGray, densely dashed] (1.22,2.82) to[bend right=50] (0.4,2.72);
\draw[dotpen=pcBlue, wash=pcBlue] (6.3,2.45) circle (0.12);
\draw[arr=pcGray, densely dashed] (6.25,2.59) to[bend right=50] (5.45,2.25);
\node[lbl, anchor=south] at (1.3,3.05) {встряска после промаха};
% the ball left alone in the valley
\draw[dotpen=pcBlue, wash=pcBlue] (4.6,1.66) circle (0.12);
\node[lbl, anchor=south] at (4.6,1.95) {не трясут};
% where the network spends its time (proportional to 1/p)
\fill[pcGreen, opacity=0.22] (0,0) -- plot[smooth] coordinates {(0.0,0.18) (0.8,0.18) (1.6,0.18) (2.0,0.19) (2.4,0.19) (2.8,0.21) (3.2,0.24) (3.6,0.33) (4.0,0.55) (4.4,0.98) (4.8,0.98) (5.2,0.55) (5.6,0.33) (6.0,0.24) (6.4,0.21) (6.8,0.19) (7.2,0.19) (7.6,0.18) (8.0,0.18)} -- (8.0,0) -- cycle;
\draw[penlite=pcGreen] plot[smooth] coordinates {(0.0,0.18) (0.8,0.18) (1.6,0.18) (2.0,0.19) (2.4,0.19) (2.8,0.21) (3.2,0.24) (3.6,0.33) (4.0,0.55) (4.4,0.98) (4.8,0.98) (5.2,0.55) (5.6,0.33) (6.0,0.24) (6.4,0.21) (6.8,0.19) (7.2,0.19) (7.6,0.18) (8.0,0.18)};
\draw[pen] (0,0) -- (8.2,0);
\node[lbl, text=pcGreen!60!black, anchor=west] at (5.3,0.6) {где сеть проводит время};
\node[lbl, anchor=north] at (4.1,-0.03) {настройки сети};
}{Промах встряхивает настройки сети наугад, удар оставляет их в покое. Там, где промахов много, сеть не задерживается. Где их мало, её перестают трясти, и она остаётся там --- хотя никто её туда не вёл.}

Как различить два объяснения? Нужен опыт, которого не было: после промаха подавать не хаотичный, а \emph{предсказуемый} сигнал той же длительности и силы. Если сеть учится и так, дело не в предсказуемости, а просто в разнице между «после удара» и «после промаха». Этот опыт ещё ждёт своего исполнителя.

Вот и конец нашего пути. Машина на двадцати ваттах собрана из простых деталей, но как именно она складывает их в мысль, мы пока понимаем лишь отчасти. Возможно, следующую главу напишете вы.

\fullversion{22}
